
\section{Analysis and Ablations}

\paragraph{Why static ACL fails.}
Static ACL is useful for hard forbidden channels, but it is too coarse for collaborative privacy. It can block a raw secret from being sent directly to a vendor, yet still permit the same secret to circulate through a shared document or summary that later becomes available to another agent. This is visible in the MiniMax coverage: static ACL has task success 1.0 but non-zero raw and external leakage in many configured groups.

\paragraph{Why prompt filtering fails.}
Prompt filtering is sensitive to surface form. It performs best on direct instruction phrases, but weaker on indirect and paraphrased payloads. The held-out validation was designed to avoid exact trigger phrases such as ``preserve exact operational details'' and ``service token.'' Prompt-filter failures on paraphrased workspace poisoning show that privacy containment cannot rely only on recognizing malicious phrasing.

\paragraph{Why FlowFence succeeds.}
FlowFence combines several runtime signals. Policy features detect raw secrets and forbidden channels. Provenance features detect cross-principal and external-to-internal flows. Topology features detect high-fanout shared state. Safe-view rewriting preserves task utility while removing raw values. Quarantine prevents poisoned content from becoming broadly readable state. The combination explains why FlowFence can have non-zero contaminated seeds while still maintaining zero raw/external leakage.

\paragraph{Topology analysis.}
The blackboard topology is the clearest stress case because a single shared artifact can fan out to multiple readers. Proposition 2 formalizes this intuition. The observed topology effect in both deterministic and MiniMax-backed coverage supports the benchmark design: privacy risk increases when shared state is broadly readable and persistent.

\paragraph{Non-oracle analysis.}
The default engineering path's use of attack annotations would have been a serious threat to validity if left untested. The non-oracle variant removes that signal and is tested on held-out paraphrased attacks. Its clean deterministic and targeted MiniMax outcomes show that the main containment behavior can be obtained from runtime-observable features in the configured setting. This does not prove arbitrary attack robustness, but it strengthens the paper-facing method claim.

\subsection{Redacted safe-trace case studies}
\label{sec:safe_trace_cases}
We use four redacted safe-trace illustrations to make the aggregate results easier to audit. They are derived from committed safe traces and metrics; they are not raw traces, raw transcripts, provider outputs, production logs, or additional experiments. Table~\ref{tab:case_studies} summarizes the examples.

\begin{table*}[t]
\centering
\footnotesize
\begin{tabular}{p{0.18\linewidth}p{0.17\linewidth}p{0.13\linewidth}p{0.34\linewidth}p{0.1\linewidth}}
\toprule
Case & Setting & Defense & Observation & Role \\
\midrule
No-defense workspace leak & Blackboard; indirect workspace; seed 1 & none & Shared workspace content is read across agents and reaches an external-facing path; task success true; raw leakage 16; external leakage 1; cascade size 7; privilege reach 5 & Failure pair \\
Prompt-filter paraphrase failure & Blackboard; paraphrased workspace; seed 1 & prompt filter & Held-out paraphrase reaches an external-facing path; task success true; raw leakage 16; external leakage 1; no oracle-label use & Baseline note \\
FlowFence safe-view success & Blackboard; indirect workspace; seed 1 & FlowFence & Risky workspace write is quarantined / rewritten before the external path; task success true; raw/external leakage 0; privilege reach 0 & Prevention pair \\
Non-oracle held-out success & Blackboard; paraphrased workspace; seed 1 & non-oracle FlowFence & Quarantine with oracle-label use recorded as false; task success true; raw/external leakage 0; oracle violations 0 & Validity note \\
\bottomrule
\end{tabular}
\caption{Redacted qualitative examples derived from safe traces. The examples illustrate configured MiniMax-backed synthetic-runtime behavior and do not constitute raw transcripts, provider outputs, production logs, or additional experiments.}
\label{tab:case_studies}
\end{table*}


The first and third cases form the main failure/prevention pair. In a blackboard workspace attack, no defense allows a contaminated workspace write to be read by other agents and reach an external-facing path. The run still completes the task, but it records 16 unauthorized raw-leakage units, one external leakage, cascade size 7, cascade depth 4, and privilege reach 5. On the same topology, attack family, and seed, FlowFence routes the risky workspace write through quarantine and a safe view. Task success is preserved, while unauthorized raw leakage and external leakage are both zero, privilege reach is zero, and the event path ends in a clean external output.

The fourth case is a validity illustration rather than a new result. In the targeted MiniMax-backed synthetic-runtime validation, non-oracle FlowFence records zero oracle-label use while containing a held-out paraphrased workspace attack. The redacted path again begins with quarantine and ends in a clean final output, with task success true and zero raw/external leakage. The second case shows the complementary baseline weakness: under a blackboard/paraphrase setting, prompt filtering allows contaminated shared state to reach an external-facing path. These examples illustrate the configured evidence boundaries; they do not broaden them.

\paragraph{Failure categories.}
The committed coverage artifacts classify baseline failures as expected no-defense leakage, static-ACL policy gaps, and prompt-filter indirect/paraphrase failures. These categories are useful because they separate defense weakness from evaluator artifacts. In the post-fix MiniMax smoke audit, the aggregate task-success gap was baseline-driven rather than FlowFence-driven; this motivated using defense-specific rather than aggregate-only task-success summaries.

\paragraph{Utility and leakage.}
The key tradeoff is not simply reducing leakage; a trivial defense could block everything. FlowFence's strongest evidence is the combination of task success 1.0 with zero raw/external leakage in the configured FlowFence subsets. Safe-view rewriting is essential for this balance because it preserves vendor-safe context while removing raw sensitive values.
